% =============================================================================
% apunte-induccion.tex
% Apunte de Induccion Matematica — UNAHUR, Licenciatura en Informatica
% Principia Computationis: Fundamentos para la Programacion
%
% Compilar con: pdflatex apunte-induccion.tex  (dos pasadas para indice)
% =============================================================================

% -----------------------------------------------------------------------------
% PREAMBULO
% -----------------------------------------------------------------------------

% Tipo de documento: article, tamano A4, fuente 11pt.
% La clase article es la adecuada para un apunte de extension media.
\documentclass[11pt,a4paper]{article}

% --- Paquetes esenciales -----------------------------------------------
\usepackage[T1]{fontenc}        % Codificacion de fuentes con caracteres latinos
\usepackage[utf8]{inputenc}     % Entrada directa de caracteres UTF-8 (tildes, enes)
\usepackage[spanish]{babel}     % Idioma espanol: rotulos, separacion silabica

% --- Matematicas -------------------------------------------------------
\usepackage{amsmath}            % Entornos matematicos: align, gather, etc.
\usepackage{amssymb}            % Simbolos adicionales: \mathbb, \square, etc.
\usepackage{amsthm}             % Teoremas, definiciones, ejemplos (entornos)
\usepackage{mathtools}          % Extension de amsmath: \coloneqq, \DeclarePairedDelimiter

% --- Diseno de pagina --------------------------------------------------
\usepackage[top=2.5cm, bottom=2.5cm, left=2.8cm, right=2.8cm]{geometry}
\usepackage{fancyhdr}           % Encabezados y pies de pagina personalizados
\usepackage{multicol}           % Texto en varias columnas (opcional)
\usepackage{enumitem}           % Control fino sobre listas enumeradas

% --- Graficos y diagramas ----------------------------------------------
\usepackage{graphicx}           % Insercion de imagenes (aunque no usemos)
\usepackage{tikz}               % Graficos vectoriales programables (diagramas)
\usepackage{xcolor}             % Colores para textos, tablas, codigo

% --- Tablas ------------------------------------------------------------
\usepackage{booktabs}           % Reglas horizontales profesionales: \toprule, \midrule
\usepackage{array}              % Control avanzado de columnas en tablas

% --- Codigo fuente -----------------------------------------------------
\usepackage{listings}           % Insercion de codigo fuente con resaltado
\lstset{
    language=C++,               % Lenguaje por defecto
    basicstyle=\ttfamily\small, % Fuente monoespaciada tamano pequenio
    keywordstyle=\bfseries\color{blue!70!black},
    commentstyle=\itshape\color{green!50!black},
    stringstyle=\color{red!70!black},
    numbers=left,               % Numeros de linea a la izquierda
    numberstyle=\tiny\color{gray},
    frame=single,               % Recuadro alrededor del codigo
    breaklines=true,            % Cortar lineas largas
    tabsize=4,
    showstringspaces=false,
    captionpos=t
}

% --- Hipervinculos (ultimo paquete en cargarse) ------------------------
\usepackage[colorlinks=true,
            linkcolor=blue!50!black,
            citecolor=green!50!black,
            urlcolor=blue!50!black,
            pdfauthor={Estudiante UNAHUR},
            pdftitle={Apunte de Induccion Matematica}]{hyperref}

% -----------------------------------------------------------------------------
% ENTORNOS DE TEOREMAS, DEFINICIONES Y EJEMPLOS
% -----------------------------------------------------------------------------
% Cada entorno cuenta con su contador ligado a la seccion actual.
% \theoremstyle{definition} usa texto normal (no italica) para definiciones.
\theoremstyle{definition}
\newtheorem{definicion}{Definici\'on}[section]

\theoremstyle{plain}
\newtheorem{teorema}{Teorema}[section]

\theoremstyle{definition}
\newtheorem{ejemplo}{Ejemplo}[section]

\theoremstyle{definition}
\newtheorem{ejercicio}{Ejercicio}[section]

\theoremstyle{remark}
\newtheorem{observacion}{Observaci\'on}[section]

% -----------------------------------------------------------------------------
% COMANDOS PERSONALIZADOS
% -----------------------------------------------------------------------------
\newcommand{\N}{\mathbb{N}}               % Numeros naturales
\newcommand{\Z}{\mathbb{Z}}               % Numeros enteros
\newcommand{\Q}{\mathbb{Q}}               % Numeros racionales
\newcommand{\R}{\mathbb{R}}               % Numeros reales
\newcommand{\ind}{{\textit{hip\'otesis inductiva}}}  % Abreviatura para HI
\newcommand{\implica}{\Rightarrow}         % Implicacion
\newcommand{\sii}{\Leftrightarrow}         % Si y solo si
\newcommand{\para}{\forall}               % Cuantificador universal
\newcommand{\existe}{\exists}             % Cuantificador existencial

% -----------------------------------------------------------------------------
% ENCABEZADOS Y PIE DE PAGINA
% -----------------------------------------------------------------------------
\pagestyle{fancy}
\fancyhf{}  % Limpia estilos anteriores
\lhead{\textbf{Inducci\'on Matem\'atica}}
\rhead{UNAHUR --- Licenciatura en Inform\'atica}
\cfoot{\thepage}
\renewcommand{\headrulewidth}{0.4pt}

% =============================================================================
% CUERPO DEL DOCUMENTO
% =============================================================================
\begin{document}

% -----------------------------------------------------------------------------
% PORTADA
% -----------------------------------------------------------------------------
\title{\Huge \textbf{Apunte de Inducci\'on Matem\'atica} \\[0.3cm]
       \large \textit{Principia Computationis} --- Fundamentos para la Programaci\'on}
\author{\textbf{Estudiante UNAHUR} \\[0.2cm]
        Licenciatura en Inform\'atica \\[0.2cm]
        Matem\'atica para Inform\'atica I / L\'ogica Computacional}
\date{\today}

\maketitle
\thispagestyle{empty}  % La portada no lleva numero de pagina
\newpage

% -----------------------------------------------------------------------------
% INDICE
% -----------------------------------------------------------------------------
\tableofcontents
\newpage

% =============================================================================
% SECCION 1: INTRODUCCION
% =============================================================================
\section{Introducci\'on}

La \textbf{inducci\'on matem\'atica} es una t\'ecnica de demostraci\'on fundamental
en matem\'atica discreta y ciencias de la computaci\'on. Permite probar que una
propiedad $P(n)$ se cumple para \textbf{todos} los n\'umeros naturales $n \in \N$,
a partir de dos ingredientes: un \emph{caso base} y un \emph{paso inductivo}.

\subsection*{Analog\'{\i}a de las fichas de domin\'o}
\label{sec:domino}

Imaginemos una hilera infinita de fichas de domin\'o colocadas en pie, una detr\'as de
otra. Queremos demostrar que, si empujamos la primera ficha, \textbf{todas} caer\'an.
Para estar seguros, necesitamos dos condiciones:
\begin{enumerate}
    \item La primera ficha cae (caso base).
    \item Si una ficha cualquiera cae, entonces la siguiente tambi\'en cae
          (paso inductivo).
\end{enumerate}
Si ambas condiciones se cumplen, podemos afirmar que toda la hilera caer\'a.
Esta es la esencia de la inducci\'on.

\subsection*{Analog\'{\i}a de los bloques numerados}

Otra imagen \'util es la de una escalera infinita de bloques numerados $1, 2, 3,
\ldots$ Queremos demostrar que podemos pararnos sobre cualquier bloque.
Si podemos pararnos sobre el bloque $1$, y desde cualquier bloque $k$ podemos
alcanzar el bloque $k+1$, entonces podemos llegar a cualquier bloque.

\subsection*{Contexto hist\'orico}

La inducci\'on matem\'atica aparece por primera vez de forma expl\'icita en el
trabajo de \textbf{Francesco Maurolico} (1575), quien la us\'o para demostrar
resultados sobre n\'umeros poligonales. Sin embargo, el nombre
\enquote{inducci\'on matem\'atica} fue acu\~nado por \textbf{Augustus De Morgan}
en 1838. La formalizaci\'on moderna proviene de los \textbf{axiomas de Peano}
(1889), en particular el axioma P5, que enunciaremos m\'as adelante.

Hoy, la inducci\'on es una herramienta cotidiana tanto en matem\'atica como en
computaci\'on: los bucles \texttt{for}, la recursi\'on y las demostraciones de
invariantes de programa se apoyan en el mismo principio l\'ogico.

% =============================================================================
% SECCION 2: EL PRINCIPIO DE INDUCCION
% =============================================================================
\section{El Principio de Inducci\'on Matem\'atica}

El principio de inducci\'on matem\'atica es una regla de inferencia que permite
demostrar proposiciones de la forma $\para n\in\N\; P(n)$.

\begin{definicion}[Principio de inducci\'on simple]
\label{def:induccion}
Sea $P(n)$ una funci\'on proposicional cuyo dominio es $\N$ (o un subconjunto
$\{n_0, n_0+1, n_0+2,\dots\}$). La demostraci\'on por inducci\'on consta de dos
pasos:
\begin{enumerate}
    \item[\textbf{Paso base:}] Se demuestra que $P(1)$ es verdadera
          (o $P(n_0)$, seg\'un el valor inicial).
    \item[\textbf{Paso inductivo:}] Se demuestra que, para todo $k\geq 1$,
          \[
          P(k) \;\implica\; P(k+1).
          \]
          La suposici\'on $P(k)$ se denomina \textbf{hip\'otesis inductiva}.
\end{enumerate}
Si ambos pasos se verifican, entonces $\para n\in\N\; P(n)$ es verdadera.
\end{definicion}

\subsection*{Regla de inferencia}

En l\'ogica proposicional, el principio de inducci\'on se escribe como:
\[
\bigl[\,P(1) \;\wedge\; \para k\;(P(k) \implica P(k+1))\,\bigr]
\;\implica\; \para n\; P(n).
\]

\subsection*{Conexi\'on con los axiomas de Peano}

El principio de inducci\'on no es un teorema que se demuestre: es un axioma.
En la axiomatizaci\'on de Peano de los n\'umeros naturales, la inducci\'on es
el quinto axioma:

\begin{definicion}[Axioma P5 de Peano]
\label{def:p5}
\[
\para P\subseteq\N:\;
\bigl[\,P(0) \;\wedge\; \para k\;(P(k) \implica P(\text{S}(k)))\,\bigr]
\;\implica\; \para n\; P(n)
\]
donde $\text{S}(k)$ es la funci\'on sucesor ($\text{S}(k)=k+1$).
\end{definicion}

Este axioma establece que \enquote{lo que vale para $0$ y se hereda, vale para
todos}. En computaci\'on, como veremos en la Secci\'on~\ref{sec:programacion},
P5 es el fundamento l\'ogico del bucle \texttt{for} y de la recursi\'on.

\subsection*{Nota sobre el valor inicial}

Aunque el axioma P5 suele formularse con $P(0)$, en muchos textos de matem\'atica
discreta (Johnsonbaugh, Rosen) se usa $P(1)$ como caso base, pues los n\'umeros
naturales se toman como $\{1,2,3,\dots\}$. La elecci\'on es convencional y no
afecta la esencia del m\'etodo.

% =============================================================================
% SECCION 3: EJEMPLOS RESUELTOS
% =============================================================================
\section{Ejemplos Resueltos}

En esta secci\'on presentamos demostraciones completas, paso a paso, para
diversas propiedades que se prueban por inducci\'on.

% --- EJEMPLO 1: Suma de 1 a n ---
\begin{ejemplo}[Suma de los primeros $n$ n\'umeros naturales]
\label{ej:suma1an}
Demostrar que para todo $n \geq 1$:
\[
1 + 2 + 3 + \cdots + n \;=\; \frac{n(n+1)}{2}.
\]

\begin{proof}
Sea $P(n)$ la proposici\'on $1+2+\cdots+n = n(n+1)/2$.

\noindent\textbf{Paso base ($n=1$):}
\[
1 = \frac{1\cdot(1+1)}{2} = \frac{2}{2} = 1,
\]
por lo que $P(1)$ es verdadera.

\noindent\textbf{Hip\'otesis inductiva:} Supongamos que $P(k)$ es verdadera para alg\'un
$k \geq 1$, es decir:
\[
1 + 2 + \cdots + k = \frac{k(k+1)}{2}.
\]

\noindent\textbf{Paso inductivo ($k \to k+1$):} Demostremos $P(k+1)$:
\[
1 + 2 + \cdots + k + (k+1).
\]
Por hip\'otesis inductiva, los primeros $k$ sumandos suman $k(k+1)/2$, luego:
\[
\begin{aligned}
1 + 2 + \cdots + k + (k+1)
&= \frac{k(k+1)}{2} + (k+1) \\
&= \frac{k(k+1) + 2(k+1)}{2} \\
&= \frac{(k+1)(k+2)}{2} \\
&= \frac{(k+1)\bigl((k+1)+1\bigr)}{2}.
\end{aligned}
\]
Esto es exactamente $P(k+1)$. Hemos demostrado $P(k) \implica P(k+1)$.

\noindent Por el principio de inducci\'on, $P(n)$ es verdadera para todo $n\geq 1$.
\qed
\end{proof}
\end{ejemplo}

% --- EJEMPLO 2: Suma de n(n+2) ---
\begin{ejemplo}[Suma de productos $n(n+2)$]
\label{ej:sumann2}
Demostrar que para todo $n \geq 1$:
\[
1\cdot 3 + 2\cdot 4 + 3\cdot 5 + \cdots + n(n+2) \;=\;
\frac{n(n+1)(2n+7)}{6}.
\]

\begin{proof}
Sea $P(n)$ la igualdad anterior.

\noindent\textbf{Paso base ($n=1$):}
\[
1\cdot 3 = 3, \qquad
\frac{1\cdot 2 \cdot (2+7)}{6} = \frac{18}{6} = 3.
\]
Se cumple $P(1)$.

\noindent\textbf{Hip\'otesis inductiva:} Para alg\'un $k\geq 1$,
\[
\sum_{i=1}^{k} i(i+2) = \frac{k(k+1)(2k+7)}{6}.
\]

\noindent\textbf{Paso inductivo ($k\to k+1$):}
\[
\begin{aligned}
\sum_{i=1}^{k+1} i(i+2)
&= \sum_{i=1}^{k} i(i+2) + (k+1)(k+3) \\
&= \frac{k(k+1)(2k+7)}{6} + (k+1)(k+3) \\
&= \frac{k(k+1)(2k+7) + 6(k+1)(k+3)}{6} \\
&= \frac{(k+1)\bigl[k(2k+7) + 6(k+3)\bigr]}{6} \\
&= \frac{(k+1)(2k^2 + 7k + 6k + 18)}{6} \\
&= \frac{(k+1)(2k^2 + 13k + 18)}{6} \\
&= \frac{(k+1)(2k^2 + 13k + 18)}{6}.
\end{aligned}
\]
Factorizamos el polinomio $2k^2+13k+18 = (2k+9)(k+2)$:
\[
\sum_{i=1}^{k+1} i(i+2)
= \frac{(k+1)(k+2)(2k+9)}{6}
= \frac{(k+1)\bigl((k+1)+1\bigr)\bigl(2(k+1)+7\bigr)}{6},
\]
que es $P(k+1)$. Luego $P(k)\implica P(k+1)$.

\noindent Por inducci\'on, $P(n)$ es verdadera para todo $n\geq 1$.
\qed
\end{proof}
\end{ejemplo}

% --- EJEMPLO 3: Desigualdad factorial ---
\begin{ejemplo}[Desigualdad factorial]
\label{ej:factorial}
Demostrar que para todo $n \geq 1$:
\[
n! \;\geq\; 2^{\,n-1}.
\]

\begin{proof}
Sea $P(n)$ la proposici\'on $n! \geq 2^{n-1}$.

\noindent\textbf{Paso base ($n=1$):}
$1! = 1$ y $2^{0} = 1$, luego $1 \geq 1$. $P(1)$ es verdadera.

\noindent\textbf{Hip\'otesis inductiva:} Supongamos $P(k)$, i.e.\ $k! \geq 2^{k-1}$.

\noindent\textbf{Paso inductivo ($k \to k+1$):}
\[
\begin{aligned}
(k+1)! &= (k+1) \cdot k! \\
      &\geq (k+1) \cdot 2^{k-1} \quad\text{(por hip\'otesis inductiva)}.
\end{aligned}
\]
Notemos que para $k \geq 1$, se cumple $k+1 \geq 2$, luego:
\[
(k+1)! \;\geq\; 2 \cdot 2^{k-1} \;=\; 2^{k}.
\]
Esto es precisamente $P(k+1)$: $(k+1)! \geq 2^{(k+1)-1} = 2^{k}$.

\noindent Por inducci\'on, $n! \geq 2^{n-1}$ para todo $n \geq 1$.
\qed
\end{proof}
\end{ejemplo}

% --- EJEMPLO 4: Divisibilidad (7^n + 1 divisible por 8) ---
\begin{ejemplo}[Divisibilidad]
\label{ej:divisibilidad}
Demostrar que $7^{n} + 1$ es divisible por $8$ para todo $n$ \textbf{impar}
positivo.

\begin{proof}
Si $n$ es impar, podemos escribir $n = 2m-1$ con $m \geq 1$.
Sea $P(m)$ la proposici\'on: $7^{2m-1} + 1$ es divisible por $8$.

\noindent\textbf{Paso base ($m=1$):}
$7^{1} + 1 = 8$, que es divisible por $8$. $P(1)$ se cumple.

\noindent\textbf{Hip\'otesis inductiva:} Supongamos $P(k)$, i.e.\
$7^{2k-1} + 1 = 8t$ para alg\'un entero $t$.

\noindent\textbf{Paso inductivo ($k \to k+1$):} Analicemos $7^{2(k+1)-1} + 1$
= $7^{2k+1} + 1$:
\[
\begin{aligned}
7^{2k+1} + 1 &= 7^{2} \cdot 7^{2k-1} + 1 \\
            &= 49 \cdot 7^{2k-1} + 1.
\end{aligned}
\]
Usamos la hip\'otesis inductiva $7^{2k-1} = 8t - 1$:
\[
\begin{aligned}
7^{2k+1} + 1 &= 49\,(8t - 1) + 1 \\
             &= 392\,t - 49 + 1 \\
             &= 392\,t - 48 \\
             &= 8\,(49\,t - 6),
\end{aligned}
\]
que es m\'ultiplo de $8$. Luego $P(k) \implica P(k+1)$.

\noindent Por inducci\'on, $7^{n}+1$ es divisible por $8$ para todo $n$ impar
positivo.
\qed
\end{proof}
\end{ejemplo}

% --- EJEMPLO 5: Funcion recursiva ---
\begin{ejemplo}[Funci\'on recursiva]
\label{ej:recursiva}
Dada la funci\'on $f: \N \to \N$ definida por:
\[
f(1) = 25, \qquad
f(n+1) = f(n) + 4,
\]
conjeturar y demostrar una f\'ormula cerrada para $f(n)$.

\begin{proof}
Calculemos los primeros valores:
\[
f(1)=25,\; f(2)=29,\; f(3)=33,\; f(4)=37,\; \dots
\]
Parece que $f(n) = 4n + 21$. Demostr\'emoslo por inducci\'on.

Sea $P(n)$ la proposici\'on $f(n) = 4n + 21$.

\noindent\textbf{Paso base ($n=1$):}
$4\cdot 1 + 21 = 25 = f(1)$. $P(1)$ es verdadera.

\noindent\textbf{Hip\'otesis inductiva:} $f(k) = 4k + 21$.

\noindent\textbf{Paso inductivo ($k \to k+1$):}
\[
\begin{aligned}
f(k+1) &= f(k) + 4 \quad\text{(definici\'on recursiva)} \\
       &= (4k + 21) + 4 \quad\text{(hip\'otesis inductiva)} \\
       &= 4k + 25 \\
       &= 4(k+1) + 21,
\end{aligned}
\]
que es $P(k+1)$.

\noindent Por inducci\'on, $f(n) = 4n + 21$ para todo $n \geq 1$.
\qed
\end{proof}
\end{ejemplo}

% =============================================================================
% SECCION 4: INDUCCION FUERTE
% =============================================================================
\section{Inducci\'on Fuerte (Forma Fuerte de Inducci\'on)}

En algunas demostraciones, la hip\'otesis inductiva $P(k)$ no es suficiente para
probar $P(k+1)$: necesitamos suponer que la propiedad vale para \textbf{todos}
los valores menores o iguales a $k$, no solo para $k$. Esto es la
\emph{inducci\'on fuerte}.

\begin{definicion}[Inducci\'on fuerte]
\label{def:induccionfuerte}
Sea $P(n)$ una funci\'on proposicional definida sobre los enteros $\geq n_0$.
Si se cumple:
\begin{enumerate}
    \item $P(n_0)$ es verdadera (paso base).
    \item Para todo $k > n_0$, la veracidad de
          $P(n_0), P(n_0+1), \dots, P(k-1)$ implica $P(k)$.
\end{enumerate}
entonces $P(n)$ es verdadera para todo $n \geq n_0$.
\end{definicion}

En forma simb\'olica:
\[
\bigl[\,P(n_0) \;\wedge\;
\para k>n_0\;
\bigl(P(n_0)\wedge P(n_0+1)\wedge\cdots\wedge P(k-1) \implica P(k)\bigr)\,\bigr]
\;\implica\; \para n\geq n_0\; P(n).
\]

\subsection*{Ejemplo: Sellos postales}

Supongamos que tenemos sellos de $2$ y $5$ centavos. Queremos demostrar que
podemos formar cualquier valor de al menos $4$ centavos.

\begin{proof}
Sea $P(n)$: \enquote{el valor $n$ puede formarse con sellos de $2$ y $5$}.
Verifiquemos los casos base:
\begin{itemize}
    \item $4 = 2+2$.
    \item $5 = 5$.
    \item $6 = 2+2+2$.
    \item $7 = 2+5$.
\end{itemize}
$P(4), P(5), P(6), P(7)$ son verdaderas.

Ahora, para $k \geq 8$, supongamos que $P(4), P(5), \dots, P(k-1)$ son
verdaderas. En particular, $P(k-4)$ es verdadera porque $k-4 \geq 4$.
Si $k-4$ puede formarse, entonces agregando un sello de $2$ obtenemos $k$.
Luego $P(k)$ es verdadera.

Por inducci\'on fuerte, todo valor $n \geq 4$ puede formarse.
\qed
\end{proof}

% =============================================================================
% SECCION 5: PROPIEDAD DEL BUEN ORDEN
% =============================================================================
\section{Propiedad del Buen Orden}

\begin{definicion}[Propiedad del buen orden]
\label{def:wellordering}
Todo conjunto no vac\'io de enteros no negativos posee un
\textbf{elemento m\'{\i}nimo}.
\end{definicion}

Esta propiedad, aparentemente obvia, es equivalente al principio de inducci\'on
matem\'atica. De hecho, se puede demostrar que:

\[
\text{Principio de inducci\'on} \;\sii\; \text{Propiedad del buen orden}.
\]

\begin{proof}[Demostraci\'on de la equivalencia (esquema)]
\noindent\textbf{($\Rightarrow$)} Supongamos el principio de inducci\'on.
Sea $A \subseteq \N$ no vac\'io. Supongamos (por contradicci\'on) que $A$ no
tiene m\'inimo. Definimos $P(n)$: \enquote{$n \notin A$}. Como $A$ no tiene
m\'inimo, $0 \notin A$ (si $0\in A$, ser\'ia el m\'inimo). Luego $P(0)$ vale.
Si $P(k)$ vale para todo $k \leq n$, entonces $n+1\notin A$ (de lo contrario
ser\'ia el m\'inimo). Por inducci\'on, $P(n)$ vale para todo $n$, luego
$A = \varnothing$, contradicci\'on. Por tanto $A$ tiene m\'inimo.

\noindent\textbf{($\Leftarrow$)} Supongamos la propiedad del buen orden.
Sea $P(n)$ tal que $P(0)$ es verdadera y $P(k) \implica P(k+1)$.
Sea $A = \{ n \in \N \mid P(n) \text{ es falsa} \}$. Si $A$ no es vac\'io,
tiene un m\'inimo $m$. Como $P(0)$ es verdadera, $m > 0$. Luego $m-1 \notin A$,
por lo que $P(m-1)$ es verdadera. Pero $P(m-1) \implica P(m)$, contradicci\'on.
Luego $A$ debe ser vac\'io y $P(n)$ vale para todo $n$.
\qed
\end{proof}

Esta equivalencia es importante porque muestra que el principio de inducci\'on
no es una \enquote{rareza} l\'ogica, sino que expresa una propiedad fundamental
de los n\'umeros naturales: todo subconjunto no vac\'io tiene un primer elemento.

% =============================================================================
% SECCION 6: INDUCCION Y PROGRAMACION
% =============================================================================
\section{Inducci\'on y Programaci\'on}
\label{sec:programacion}

La inducci\'on matem\'atica no es solo una herramienta te\'orica: es el
fundamento l\'ogico de dos conceptos computacionales centrales:
los bucles y la recursi\'on.

\subsection{Bucle \texttt{for} como inducci\'on}

Un bucle \texttt{for} recorre los valores $i = 1, 2, \dots, n$ y mantiene una
propiedad que se cumple al inicio de cada iteraci\'on: el
\textbf{invariante de bucle}. La demostraci\'on de que el invariante se
mantiene es exactamente el paso inductivo, y la verificaci\'on inicial es el
caso base.

\begin{ejemplo}[Suma con \texttt{for}]
\label{ej:for}
El siguiente programa calcula $1+2+\cdots+n$:

\begin{lstlisting}[language=C++]
// Calcula la suma de 1 a n usando un bucle for.
// Invariante: al inicio de cada iteracion i,
//   total almacena 1 + 2 + ... + (i-1).
int sumar(int n) {
    int total = 0;           // caso base: suma de 0 términos = 0
    for (int i = 1; i <= n; i++) {
        total += i;          // paso inductivo: total += i
    }
    return total;            // post: total = n(n+1)/2
}
\end{lstlisting}

La correspondencia es directa:
\begin{center}
\begin{tabular}{lll}
\toprule
\textbf{Etapa} & \textbf{Inducci\'on} & \textbf{Bucle \texttt{for}} \\
\midrule
Inicio & $P(0)$: suma de 0 t\'erminos & \texttt{total = 0} \\
Iteraci\'on $i$ & $P(i-1)$ es verdadera & Invariante se cumple \\
Cuerpo & Demostramos $P(i-1)\implica P(i)$ & \texttt{total += i} \\
Final & $\para i\; P(i)$ & \texttt{total} contiene la suma \\
\bottomrule
\end{tabular}
\end{center}
\end{ejemplo}

\subsection{Recursi\'on como inducci\'on estructural}

Una funci\'on recursiva se define con uno o m\'as casos base y una llamada
recursiva con argumentos m\'as pequenios. Demostrar que la funci\'on es correcta
equivale a demostrar:
\begin{itemize}
    \item El caso base produce el resultado correcto (paso base).
    \item Si la funci\'on funciona para argumentos menores, funciona para el
          argumento actual (paso inductivo).
\end{itemize}

\begin{ejemplo}[Funci\'on factorial recursiva]
\label{ej:recursiva2}

\begin{lstlisting}[language=C++]
// Factorial recursivo: n! = n * (n-1)!
// Caso base: 0! = 1
// Paso inductivo: n! = n * (n-1)!
int factorial(int n) {
    if (n == 0) return 1;              // caso base
    return n * factorial(n - 1);       // paso inductivo
}
\end{lstlisting}
\end{ejemplo}

\subsection{Invariante de bucle}

El invariante de bucle es una propiedad que:
\begin{enumerate}
    \item Es verdadera \textbf{antes} de la primera iteraci\'on (caso base).
    \item Si es verdadera al inicio de una iteraci\'on, sigue siendo verdadera
          al final de esa iteraci\'on (paso inductivo).
    \item Al terminar el bucle, el invariante nos da informaci\'on \'util sobre
          el resultado.
\end{enumerate}
Esta t\'ecnica de demostraci\'on de programas fue formalizada por
\textbf{C. A. R. Hoare} (1969) y es la aplicaci\'on m\'as directa de la
inducci\'on en ingenier\'{\i}a de software.

% =============================================================================
% SECCION 7: EJERCICIOS PROPUESTOS
% =============================================================================
\section{Ejercicios Propuestos}

Los siguientes ejercicios se dejan como pr\'actica para el estudiante.
Se recomienda intentar cada demostraci\'on por inducci\'on antes de consultar
una soluci\'on.

\begin{ejercicio}[Suma de los primeros $n$ impares]
\label{ej:impares}
Demostrar que la suma de los primeros $n$ n\'umeros impares es $n^{2}$:
\[
1 + 3 + 5 + \cdots + (2n-1) = n^{2}.
\]
\end{ejercicio}

\begin{ejercicio}[Suma de los primeros $n$ cubos]
\label{ej:cubos}
Demostrar que:
\[
1^{3} + 2^{3} + 3^{3} + \cdots + n^{3} =
\left[\frac{n(n+1)}{2}\right]^{2}.
\]
\end{ejercicio}

\begin{ejercicio}[Suma de $n(n+1)$]
\label{ej:sumann1}
Demostrar que para todo $n \geq 1$:
\[
1\cdot 2 + 2\cdot 3 + 3\cdot 4 + \cdots + n(n+1) =
\frac{n(n+1)(n+2)}{3}.
\]
\end{ejercicio}

\begin{ejercicio}[Diagonales de un pol\'{\i}gono]
\label{ej:diagonales}
Demostrar que un pol\'{\i}gono convexo de $n \geq 3$ lados tiene
\[
d_n = \frac{n(n-3)}{2}
\]
diagonales.
\end{ejercicio}

\begin{ejercicio}[Desigualdad $2^{n} \geq n^{2}$]
\label{ej:desigualdad2n}
Demostrar que $2^{n} \geq n^{2}$ para todo $n \geq 4$.
\end{ejercicio}

\begin{ejercicio}[Suma de Fibonacci]
\label{ej:fibonacci}
Sea $F_n$ la sucesi\'on de Fibonacci definida por:
\[
F_1 = 1,\; F_2 = 1,\; F_n = F_{n-1} + F_{n-2}\;(n\geq 3).
\]
Demostrar que la suma de los primeros $n$ t\'erminos es:
\[
F_1 + F_2 + \cdots + F_n = F_{n+2} - 1.
\]
\end{ejercicio}

\begin{ejercicio}[Suma de cuadrados]
\label{ej:cuadrados}
Demostrar que:
\[
1^{2} + 2^{2} + 3^{2} + \cdots + n^{2} = \frac{n(n+1)(2n+1)}{6}.
\]
\end{ejercicio}

\begin{ejercicio}[Suma de potencias de $2$]
\label{ej:potencias2}
Demostrar que:
\[
2^{0} + 2^{1} + 2^{2} + \cdots + 2^{n} = 2^{n+1} - 1.
\]
\end{ejercicio}

% =============================================================================
% SECCION 8: FALACIAS Y ERRORES COMUNES
% =============================================================================
\section{Falacias y Errores Comunes}

La inducci\'on matem\'atica es una herramienta poderosa, pero debe aplicarse
correctamente. A continuaci\'on presentamos un ejemplo cl\'asico de falacia
para ilustrar lo que puede salir mal.

\subsection*{La paradoja de los caballos}

\enquote{Todos los caballos son del mismo color}. Intentemos \enquote{demostrarlo}
por inducci\'on.

Sea $P(n)$: \enquote{en cualquier conjunto de $n$ caballos, todos son del mismo
color}.

\noindent\textbf{Paso base ($n=1$):} En un conjunto de un solo caballo,
trivialmente todos son del mismo color. $P(1)$ es verdadera.

\noindent\textbf{Hip\'otesis inductiva:} Supongamos $P(k)$: en cualquier
conjunto de $k$ caballos, todos son del mismo color.

\noindent\textbf{Paso inductivo ($k\to k+1$):} Tomemos un conjunto de $k+1$
caballos. Numer\'emoslos $c_1, c_2, \dots, c_{k+1}$.
\begin{itemize}
    \item El subconjunto $\{c_1, \dots, c_k\}$ tiene $k$ caballos, luego por
          $P(k)$ todos son del mismo color.
    \item El subconjunto $\{c_2, \dots, c_{k+1}\}$ tambi\'en tiene $k$ caballos,
          luego todos son del mismo color.
    \item Los caballos $c_2, \dots, c_k$ est\'an en ambos subconjuntos, por lo
          que el color de $c_1$ es el mismo que el de $c_2$, que es el mismo
          que el de $c_{k+1}$. As\'{\i}, todos los $k+1$ caballos son del mismo
          color.
\end{itemize}

\noindent\textbf{Conclusi\'on:} Todos los caballos son del mismo color.

\begin{observacion}[\texorpdfstring{¿D\'onde falla la demostraci\'on?}{Donde falla la demostracion}]
\label{obs:falla}
El error est\'a en el paso inductivo para $k=1$. Cuando tenemos $k+1 = 2$
caballos, los subconjuntos $\{c_1\}$ y $\{c_2\}$ tienen cada uno $1$ caballo,
pero no hay caballos en com\'un entre ellos. Por tanto, no podemos inferir que
$c_1$ y $c_2$ tengan el mismo color. La hip\'otesis inductiva $P(1)$ es
verdadera, pero $P(1) \implica P(2)$ es falso.

La lecci\'on: siempre hay que verificar que el paso inductivo funcione para
\textbf{todos} los valores de $k$, especialmente los m\'as pequenios.
\end{observacion}

\subsection*{Otros errores comunes}

\begin{itemize}
    \item \textbf{Olvidar el paso base:} Sin caso base, la inducci\'on no tiene
          punto de partida. Podr\'{\i}amos \enquote{demostrar} cualquier cosa.
    \item \textbf{Paso base incompleto:} A veces el paso inductivo requiere
          m\'as de un valor base (como en la inducci\'on fuerte).
    \item \textbf{Hip\'otesis inductiva mal formulada:} La hip\'otesis debe
          ser exactamente $P(k)$, no una versi\'on modificada.
    \item \textbf{Variable cambiada:} No confundir la variable $n$ con $k$.
          $k$ es un valor fijo arbitrario; $n$ es la variable universal.
\end{itemize}

% =============================================================================
% SECCION 9: BIBLIOGRAFIA
% =============================================================================
\section{Bibliograf\'{\i}a}

\begin{enumerate}[label={[}\arabic{*}{]}]
    \item Johnsonbaugh, R. \textit{Matem\'aticas Discretas}. 6.\textsuperscript{a}
          ed., Pearson, 2005. — Cap\'{\i}tulos 1.7 y 1.8 (Inducci\'on Matem\'atica
          y Forma Fuerte de Inducci\'on).
    \item Rosen, K. \textit{Discrete Mathematics and its Applications}.
          8.\textsuperscript{a} ed., McGraw-Hill, 2019. — Cap\'{\i}tulo 5
          (Inducci\'on y Recursi\'on).
    \item Sominsky, I. \textit{El M\'etodo de Inducci\'on Matem\'atica}.
          Editorial MIR, Mosc\'u, 1975. — Monograf\'{\i}a cl\'asica dedicada
          \'unicamente a la inducci\'on.
    \item Peano, G. \textit{Arithmetices principia, nova methodo exposita}.
          1889. — Formulaci\'on original de los axiomas de Peano, incluido
          el axioma P5 de inducci\'on.
    \item Hoare, C. A. R. \enquote{An Axiomatic Basis for Computer Programming}.
          \textit{Communications of the ACM}, vol.\ 12, n.\ 10, 1969,
          pp.\ 576--580.
    \item Velleman, D. \textit{How to Prove It: A Structured Approach}.
          3.\textsuperscript{a} ed., Cambridge University Press, 2019.
    \item Material de c\'atedra: \texttt{REFERENCIA\_AXIOMAS\_ZFC\_PEANO.md}
          (Principia Computationis: axiomas de Peano, ZFC y su relaci\'on
          con C++).
\end{enumerate}

% -----------------------------------------------------------------------------
% DIAGRAMA TIKZ: FICHAS DE DOMINO (ilustracion opcional)
% -----------------------------------------------------------------------------
\section*{Ap\'endice: Diagrama de las fichas de domin\'o}

Como ilustraci\'on de la analog\'{\i}a de la Secci\'on~\ref{sec:domino},
incluimos un diagrama TikZ simple que muestra las primeras fichas.

\vspace{0.5cm}

\begin{center}
\begin{tikzpicture}[scale=0.9]
    % Dibujamos 6 fichas de domino cayendo en secuencia
    \foreach \x in {0,1,2,3,4,5} {
        \draw[fill=gray!20] (\x*1.8,0) rectangle (\x*1.8+1.2,2);
        \draw (\x*1.8+0.6,1) -- (\x*1.8+0.6,1.5);
        \fill[black] (\x*1.8+0.4,1.3) circle (0.08);
        \fill[black] (\x*1.8+0.8,0.7) circle (0.08);
        % Flecha de caida
        \draw[->,thick,red!70!black] (\x*1.8+0.6,2.5) -- (\x*1.8+0.6,2.1);
        % Etiqueta debajo
        \node[below] at (\x*1.8+0.6, -0.2) {\small $P(\x+1)$};
    }
    % Etiqueta general
    \node[above] at (5.4, 3.2) {\textbf{Ca\'{\i}da de fichas de domin\'o}};
    \node[below] at (5.4, -0.8) {Si $P(1)$ cae y $P(k)\implica P(k+1)$, entonces todas caen.};
\end{tikzpicture}
\end{center}

\vspace{0.3cm}

\begin{center}
\begin{tikzpicture}[scale=0.9]
    % Escalera de la induccion: bloques numerados
    \foreach \x in {0,...,4} {
        \draw[fill=blue!15] (\x*1.5,\x*0.5) rectangle (\x*1.5+1,\x*0.5+0.6);
        \node at (\x*1.5+0.5,\x*0.5+0.3) {$\x+1$};
        \ifnum\x<4
            \draw[->,thick] (\x*1.5+1,\x*0.5+0.3) -- (\x*1.5+1.5,\x*0.5+0.8);
        \fi
    }
    \node[above] at (3.75, 2.8) {\textbf{Escalera de la inducci\'on}};
    \node[below] at (3.75, -0.3) {De cada bloque se alcanza el siguiente.};
\end{tikzpicture}
\end{center}

% =============================================================================
% FIN DEL DOCUMENTO
% =============================================================================
\end{document}
